\section{Research Gap and Related Work}
\label{sec:related}

\paragraph{Search method.}
We positioned SAT-SA with a structured search rather than a systematic review. On
27 September 2026 we ran 20 queries covering SOC assessment and metrics, alert
prioritisation and alert fatigue, security process mining and conformance checking,
peer-group analytics, supervisory technology, and tamper-evident and post-quantum
audit records. Each query was run against OpenAlex (title and abstract) and Crossref,
and the 25 top-ranked records per query and source were retrieved. This gave 656
unique records after deduplication. One screener read every title and, for 106
records, the abstract; 51 records were included and joined with the 26 sources of an
earlier focused review. arXiv throttled the automated queries and was not searched,
and records beyond rank 25 were not screened. The absence of a work in this search
is therefore not evidence that it does not exist. Queries, counts, screening
decisions with reasons, and a comparison matrix are in the supplementary material.

\paragraph{SOC assessment and performance measurement.}
SOC-CMM scores capability maturity across business, people, process, technology and
services domains through a self-assessment tool~\cite{vanos2016soccmm}, and
CTI-SOC2M2 uses the degree of threat-intelligence integration as a proxy for SOC
service maturity~\cite{schlette2021ctisoc2m2}. Agyepong et al.\ review SOC analyst
challenges and metrics~\cite{agyepong2019review} and weight analyst-performance
indicators using a Delphi panel and the analytic hierarchy
process~\cite{agyepong2023analyst}; Forsberg and Frantti derive technical SOC
metrics through design science~\cite{forsberg2023metrics}. A case study in one
academic SOC found that analysts who scored well on key performance indicators did
not necessarily further the SOC's goals~\cite{morin2025mismatched}, which cautions
against reading operational metrics as quality. SAIBERSOC evaluates the performance
of operational SOCs empirically by injecting synthetic attacks and measuring outputs
such as detection accuracy and time to investigation, with 124 students acting as
analysts~\cite{rosso2020saibersoc}. Surveys and qualitative studies characterise SOC
building blocks, challenges and human
factors~\cite{vielberth2020soc,kokulu2019soc,sundaramurthy2015burnout,sundaramurthy2016contradictions}.
These instruments measure maturity, define metrics, explain practice or test a SOC
by injection. None of them computes findings from records that a monitored SOC
submits to an external supervisor.

\paragraph{Record-based assessment of incident handling.}
The closest prior work is the compliance assessment system of Palma et
al.~\cite{palma2024compliance}. It aligns an organisation's incident-management
event log with a reference process derived from frameworks such as ITIL and ISO
27035, classifies deviations with a taxonomy, prioritises them with a cost model,
and supports an auditor. It was validated on a benchmark and demonstrated on a public
real incident-management log. A companion benchmark compares such process-based
assessment models quantitatively~\cite{palma2024benchimp}. Automated, record-based
assessment of how incidents were handled, aimed at an auditor, therefore already
exists. Earlier work applied process mining to IT incident management against
ITIL~\cite{ferreira2008itil} and to security audit trails~\cite{vanderaalst2005security}.
Recent work represents incident-response activities and responder communications
from tool logs as a knowledge graph for process analysis~\cite{galadima2026irgraph},
and DomainPrio analyses the population of a SOC's investigations rather than single
alerts~\cite{chiba2022domainprio}. SAT-SA differs from Palma et al.\ in mechanism
and scope rather than in kind. It uses rule-encoded checks over SOC evidence
categories (alerts, cases, investigation steps, escalations, dispositions, assets)
instead of alignment with a process model. It compares an entity with a declared peer
cohort, ranks monitored organisations by a fused risk, and binds the supervisor's
decision to the analysed state cryptographically. We did not compare SAT-SA with
that system experimentally.

\paragraph{Alert prioritisation and alert fatigue.}
A large literature ranks alerts or incidents for SOC analysts. Jalalvand et
al.\ review 89 prioritisation papers through a human--AI teaming
lens~\cite{jalalvand2024prioritisation}, and Tariq et al.\ survey alert-fatigue
mitigations~\cite{tariq2025alertfatigue}. Systems include resource-aware
optimisation of which alerts to investigate~\cite{shah2019alertselection},
provenance-based triage (NoDoze~\cite{hassan2019nodoze}), guided response trained on
customer triage labels~\cite{freitas2024guide}, large-scale incident
ranking~\cite{freitas2026aip} and learning to defer to analysts~\cite{jalalvand2025l2d}.
SAT-SA differs in unit and user: it ranks \emph{monitored organisations} for a
\emph{supervisor} from rule-based findings about how work was handled, not individual
alerts for an analyst. It is not a competitor to these systems and is not compared
with them.

\paragraph{Process analytics and absence detection.}
Process mining discovers and checks processes from event
logs~\cite{vanderaalst2016pm}. Conformance checking relates recorded behaviour to a
model and exposes skipped or out-of-order
activities~\cite{rozinat2008conformance,carmona2018conformance}. SAT-SA's
execution-gap and ``negative-space'' workers are rule-encoded conformance checks (for
example, an acknowledged alert whose case has no investigation step). We treat them
as an adaptation of this prior art, not a new technique.

\paragraph{Peer comparison and robust deviation.}
Peer group analysis flags entities that diverge from similar entities over
time~\cite{weston2008pga}, in the broader statistical fraud-detection
tradition~\cite{bolton2002fraud}. SAT-SA's peer worker compares an entity with a
declared cohort using the median and the median absolute deviation
(MAD)~\cite{leys2013mad}, with a minimum cohort size and tenant isolation.

\paragraph{Supervisory technology and human oversight.}
Financial supervisors increasingly use technology to analyse regulated entities'
data (``suptech'')~\cite{broeders2018suptech}, and human supervisory judgement
remains indispensable there, as in SAT-SA's design. Automation bias in human use of
decision support~\cite{parasuraman2010automation} has been reviewed specifically for
SOCs~\cite{tilbury2024automationbias}; it motivates keeping the decision with the
supervisor, and it was not studied here.

\paragraph{Integrity and provenance of records.}
Forward-secure and tamper-evident logs detect alteration of recorded
history~\cite{schneier1999logs,crosby2009logging}. Practical systems protect
operating-system audit logs with trusted execution~\cite{paccagnella2020custos} or
combine ratcheting with keystream storage~\cite{guardiola2022sealfs}, and recent work
applies post-quantum primitives to tamper-evident forensic logging in the
cloud~\cite{fugkeaw2026pqabl}. W3C PROV provides a provenance
vocabulary~\cite{w3c2013provdm}. SAT-SA does not propose a logging scheme or new
cryptography. It composes standard primitives, SHA3-256~\cite{nist2015fips202} and
ML-DSA-65~\cite{nist2024fips204}, to bind one canonical supervisory record to the
authorised decision.

\paragraph{Gap, stated conservatively.}
Record-based assessment of incident handling~\cite{palma2024compliance}, supervisory
analysis of submitted data~\cite{broeders2018suptech}, peer comparison and
tamper-evident records each exist. In the structured search we did not find a work
that combines them in the way SAT-SA does: supervisor-side analysis of SOC evidence
submissions across monitored organisations, with record-cited findings, peer-cohort
comparison, entity risk ranking and a cryptographically bound human decision. This
is a combination of established parts, identified by a search with the limits
stated above. We therefore offer the framing only as a low-confidence candidate
contribution and make no novelty claim.
